\section{From atom types to hereditarily finite sets}
\label{sec:hf-transfer}

The preceding counting equivalence concerns tuples of input atoms.  A
choiceless computation also constructs sets of atoms, sets of such sets, and
objects of unbounded finite rank.  We now show how small supports let us
transfer bounded-variable counting formulas on a domain containing all these
objects.  The domain itself will be countably infinite; every object in it
is nevertheless hereditarily finite.  The counting argument below respects
this distinction. The support, molecule and form method originates in
Blass, Gurevich and Shelah
\cite[Sections~8--9 of the published version]{bgs1999} and was developed
for counting by Dawar, Richerby and Rossman
\cite[Section~8 of the author version]{drr2008}. We give the full
quantitative argument for the chosen automorphism subgroups and support
bound used here, including exact counts of represented values.

We first isolate the precise hypotheses needed from the preceding sections.
Write $\mathrm{HF}(A)$ for the atoms of $A$ together with all
hereditarily finite sets built from them.
For a finite relational structure $A$, let $\operatorname{tp}_M^A(\bar a)$
be the set of all $C^M$ formulas true at $\bar a$, with free variables among
the indicated tuple positions.  Tuples may have repeated entries.  Two
such types are compared using the same ordered list of variable positions.

\begin{theorem}[Transfer through small supports]
\label{thm:hf-transfer}
Let $A_0,A_1$ be nonempty finite structures in the same binary relational
vocabulary, and let $K_i\leq\operatorname{Aut}(A_i)$ for $i\in\{0,1\}$.
Let $s,m,M$ be positive integers satisfying
\[
 \max\{4,m+1\}s\leq M.
\]
Assume the following two properties.
\begin{enumerate}
\item $A_0$ and $A_1$ satisfy the same $C^M$ sentences.
\item In each $A_i$, any two tuples of length at most
      $\max\{2,m\}s$ with the same $C^M$ type lie in the same $K_i$ orbit.
\end{enumerate}
For $i\in\{0,1\}$, let $D_i$ consist of the objects $x$ in
$\mathrm{HF}(A_i)$ such that $x$ and every object in its transitive closure
have a $K_i$-support of length at most $s$.  Equip $D_i$ with membership,
an atom predicate, and the relations of $A_i$, interpreted as false on
arguments containing a non-atom.  Then $D_0$ and $D_1$ satisfy the same
$C^m$ sentences.
\end{theorem}

Here atoms are urelements, distinct from all sets.  A tuple supports an
object if its pointwise stabilizer fixes that object.  The action of $K_i$
on atoms extends recursively to sets.  We can pad every support to exactly
$s$ coordinates, using repetitions or arbitrary atoms, because $A_i$ is
nonempty.  The resulting domains $D_i$ contain all atoms and all pure
hereditarily finite sets.  They are transitive under membership, countable,
and invariant under $K_i$: if $\alpha$ supports $x$, then $g\alpha$
supports $gx$.

For our application, $A_i$ in Theorem~\ref{thm:hf-transfer} is the analytic
expansion $\bar A_{b^i}$, and $K_i$ is its central flow group $K$.  The base
counting-equivalence and homogeneity results provide its two assumptions
with the parameter bounds established at the end of
Section~\ref{sec:base-counting}. For the remainder of the paper, write
$\mathcal D_{b^i}$ for this supported
domain over $\bar A_{b^i}$, using the support length in
\eqref{eq:support-length}. The atom predicate is definable in our domain by
$\mathsf{Ed}(x,x)\lor\mathsf{Cf}(x,x)$.  No analytic expansion is supplied
to the program; it is used only to prove the indistinguishability theorem.

The proof has three stages: count extensions of atom types; represent
supported objects and transfer equality and membership; then match the
whole domains while retaining the assigned parameters.

\subsection{Exact extension counts for atom types}

All types in the next lemma are realized in $A_0$ or $A_1$.  Thus there are
only finitely many types of any fixed tuple length, even though a type
records infinitely many formulas.

\begin{lemma}[Extension counts]
\label{lem:hf-type-counts}
Suppose $P$ and $P'$ have the same $C^M$ type in two structures chosen,
with repetition allowed, from $A_0,A_1$.  If $|P|+d\leq M$, then for every
realized type $\rho$ of length $|P|+d$, the numbers of tuples $\gamma$
and $\gamma'$ satisfying
\[
 \operatorname{tp}_M(P,\gamma)=\rho,
 \qquad
 \operatorname{tp}_M(P',\gamma')=\rho
\]
are equal.  In particular a type realized over one prefix is realized over
the other.  Equal types also have equal types on every projected or
reordered subtuple.
\end{lemma}

\begin{proof}
Projection and reordering follow by renaming the free tuple variables:
an injection between the selected position sets extends to a permutation
of the $M$ available variable names.  Repeated atom values cause no
problem; positions remain distinct, and equality formulas record all
coincidences of values.

First let $d=1$ and write $a=|P|$.  A realized type $\rho$ of $(a+1)$-tuples
can be isolated among the tuples of both finite structures by a single
$C^M$ formula $\theta_\rho$.  Indeed, choose a representative of $\rho$.
For every tuple in the two structures whose type differs from $\rho$,
choose a formula distinguishing it from the representative, negating the
formula if necessary.  The conjunction of this finite collection is true
exactly on the tuples of type $\rho$.  Its free variables lie among the
first $a+1$ positions, and it still uses only the same $M$ variable names.
If all tuples have type $\rho$, use a tautology.

If exactly $h$ extensions of $P$ realize $\rho$, the formula
\[
 \exists^{=h}x_{a+1}\,
       \theta_\rho(x_1,\ldots,x_a,x_{a+1})
\]
is true at $P$.  Equality of prefix types makes it true at $P'$ as well.
This proves the one-coordinate assertion, including $h=0$ and $a=0$.
For empty prefixes in different structures, the required equality of
types is precisely their agreement on $C^M$ sentences.

For $d>1$, project $\rho$ onto the first $a+d-1$ positions, obtaining
$\sigma$.  By induction the number of extensions of $P$ to type $\sigma$
equals the corresponding number over $P'$.  The one-coordinate argument,
used both within a structure and between the two structures, shows that
every tuple of type $\sigma$ has the same number of last-coordinate
extensions to type $\rho$.  Multiplying these two finite counts proves
the assertion.
\end{proof}

\subsection{A common syntax for supported objects}

A \emph{molecule} is a tuple $\alpha\in A_i^s$.  It records a support, but
need not be a minimal support.  We describe sets by a common collection
of formal expressions, independent of the side $i$.

There are $s$ atomic forms $c_1,\ldots,c_s$.  A set form is a finite set
of pairs $(\psi,\rho)$, where $\psi$ is a previously constructed form
and $\rho$ is a realized $C^M$ type of length $2s$ in $A_0$ or $A_1$.
These are syntactic constructors: atomic forms are distinct from every
set form, including the empty set form.  A form has finite depth; set
forms have greater depth than their constituent forms.

For a molecule $\alpha\in A_i^s$, define the value of a form recursively:
\begin{align}
 c_j*\alpha&=\alpha_j, \label{eq:hf-atomic-form}\\
 \phi*\alpha&=
 \{\psi*\beta:
        (\psi,\rho)\in\phi,
        \ \beta\in A_i^s,
        \ \operatorname{tp}_M^{A_i}(\beta,\alpha)=\rho\}
       \quad\text{if $\phi$ is a set form}.
       \label{eq:hf-set-form}
\end{align}
This produces a finite set because both the form and the base universe
are finite.  The empty set form evaluates to the empty set.

\begin{lemma}[Representation]
\label{lem:hf-representation}
Every value $\phi*\alpha$ belongs to $D_i$ and is supported by $\alpha$.
Conversely, every object of $D_i$ has such a representation.  Moreover,
for every $g\in K_i$,
\[
 g(\phi*\alpha)=\phi*(g\alpha).
\]
\end{lemma}

\begin{proof}
Automorphisms preserve $C^M$ types.  Induction on the form therefore gives
the displayed equivariance identity: in the set case, change the witness
molecule $\beta$ to $g\beta$ in~\eqref{eq:hf-set-form}.  If $g$ fixes
$\alpha$, this identity says that it fixes $\phi*\alpha$.  Induction also
puts every member of a set-form value, and all of that member's recursive
constituents, in $D_i$.  Atomic values have the required support directly.

For the converse, induct on the hereditary rank of the object.  An atom
has the form $c_j*\alpha$ for a molecule containing it.  Let $x$ be a set
in $D_i$, and choose a supporting molecule $\alpha$.  For each $y\in x$,
the induction hypothesis gives $y=\psi_y*\beta_y$.  Form the finite set
\[
 \phi=\{(\psi_y,\operatorname{tp}_M(\beta_y,\alpha)):y\in x\}.
\]
Every member of $x$ lies in $\phi*\alpha$.  Conversely, a member of
$\phi*\alpha$ has the form $\psi_y*\beta$, where
$(\beta,\alpha)$ and $(\beta_y,\alpha)$ have the same type.  Homogeneity
for tuples of length $2s$ supplies an element $g\in K_i$ mapping the
second tuple to the first.  In particular $g$ fixes $\alpha$, so it fixes
$x$.  Equivariance gives $\psi_y*\beta=gy\in x$.  Thus
$x=\phi*\alpha$.  For $x=\varnothing$, the construction gives the empty
set form and the same conclusion.
\end{proof}

Representation reduces arbitrary finite rank to finite syntax.  We next
show that the syntax can be evaluated on either base structure without
changing equality or membership, provided the supporting molecules have
the same joint type.

\begin{lemma}[Transfer of atomic facts]
\label{lem:hf-atomic-transfer}
Let $\phi,\psi$ be forms.  Suppose $(\alpha,\beta)$ and
$(\alpha',\beta')$ have the same $C^M$ type, where the two tuples may be
in the same base structure or in different base structures.  Then
\begin{align*}
 \phi*\alpha=\psi*\beta
   &\quad\Longleftrightarrow\quad
       \phi*\alpha'=\psi*\beta',\\
 \phi*\alpha\in\psi*\beta
   &\quad\Longleftrightarrow\quad
       \phi*\alpha'\in\psi*\beta'.
\end{align*}
The atom predicate and all base relations also agree on corresponding
form values.
\end{lemma}

\begin{proof}
For equality, induct on the sum of the depths of the two forms.  For two
atomic forms, equality is an equality of molecule coordinates and is
recorded by their joint type.  An atomic value never equals a set-form
value.  It remains to compare two set forms.

Suppose $\phi*\alpha=\psi*\beta$.  Take a member of
$\phi*\alpha'$ and a witness for it, say $(\theta,\rho)\in\phi$ and
$\gamma'$ with $\operatorname{tp}_M(\gamma',\alpha')=\rho$.
Lemma~\ref{lem:hf-type-counts} extends the matching prefix
$(\alpha,\beta)$ to a tuple $(\alpha,\beta,\gamma)$ with the type of
$(\alpha',\beta',\gamma')$.  Hence $\theta*\gamma$ is a member of
$\phi*\alpha$, and therefore of $\psi*\beta$.  Choose a witness
$(\xi,\sigma)\in\psi$ and a molecule $\zeta$ such that
\[
 \operatorname{tp}_M(\zeta,\beta)=\sigma,
 \qquad \theta*\gamma=\xi*\zeta.
\]
Apply the extension lemma once more, this time to the primed prefix of
length $3s$, to obtain $\zeta'$ such that
\[
 \operatorname{tp}_M(\alpha,\beta,\gamma,\zeta)
   =\operatorname{tp}_M(\alpha',\beta',\gamma',\zeta').
\]
This is the step requiring $4s\leq M$.  Projection preserves the two
member-witness types.  The induction hypothesis applies to $\theta,\xi$,
whose depths are strictly smaller than those of $\phi,\psi$, and gives
$\theta*\gamma'=\xi*\zeta'$.  Thus the chosen member belongs to
$\psi*\beta'$.  Interchanging $\phi$ and $\psi$ proves the reverse
inclusion.  Finally interchange primed and unprimed data to prove both
directions of equality transfer.  The argument includes empty set forms:
when there is no member to choose, the relevant inclusion is vacuous.

For membership, the right-hand value is either an atom, in which case
membership is false on both sides, or the value of a set form.  In the
latter case, a member witness $\theta*\gamma$ equal to $\phi*\alpha$
can be transferred over the prefix $(\alpha,\beta)$ using $3s$ positions.
Equality transfer, already established for all forms, then transfers the
witness.  Interchanging sides proves the converse.

A form is atomic on both sides or a set form on both sides.  Consequently
the atom predicate transfers.  A base relation is false if any argument
is a set; otherwise it is a relation on selected molecule coordinates,
and its truth is part of their joint base type.  Repeated arguments are
allowed throughout.
\end{proof}

\subsection{Matching entire domains while retaining parameters}

An object may have several representations $\phi*\alpha$. Thus equal
counts of molecules do not by themselves give equal counts of objects.
We will group represented values into stabilizer orbits and account for
the number of molecules representing each value. This will give a
bijection between the two domains after fixing the values of at most
$m-1$ variables. The construction is a mathematical comparison; it is
not required to be definable or executable by a choiceless program.

Call paired assignments \emph{represented alike} if their values have
representations
\[
 x_j=\phi_j*\alpha_j,\qquad x'_j=\phi_j*\alpha'_j
 \quad(1\leq j\leq r),
\]
using common forms, and the concatenated tuples
$P=(\alpha_1,\ldots,\alpha_r)$ and
$P'=(\alpha'_1,\ldots,\alpha'_r)$ have the same $C^M$ type.
Dropping a position preserves this condition by projection.  Empty
assignments are represented alike by the sentence-equivalence assumption.

\begin{lemma}[A bijection extending retained assignments]
\label{lem:hf-bijection}
For any represented-alike assignments with $r\leq m-1$, there is a
bijection $f:D_0\longrightarrow D_1$ such that adjoining any pair
$x,f(x)$ gives represented-alike assignments.  In particular $f(x_j)=x'_j$
for every retained position.
\end{lemma}

\begin{proof}
Fix representing tuples $P,P'$.  A label is a pair $(\phi,\tau)$,
where $\phi$ is a form and $\tau$ is a realized base type of length
$(r+1)s$.  Associate to it the finite value sets
\[
 X_{\phi,\tau}(P)=
   \{\phi*\gamma:\gamma\in A_0^s,
             \ \operatorname{tp}_M(P,\gamma)=\tau\}
\]
and $X_{\phi,\tau}(P')$ on the other side.  The extension-count lemma
shows that the two sets are either both empty or both nonempty.

Every nonempty value set is a single orbit of the stabilizer $K_{0,P}$.
Indeed, any two molecules with the specified type over $P$ are carried
to one another by an element fixing $P$, using homogeneity for
$(r+1)s\leq ms$ positions.  Equivariance carries their values along with
them.  Conversely, an element fixing $P$ preserves the specified joint
type, and thus preserves the value set.  The same assertion holds on the
other side.  It follows that two nonempty label sets on one side either
coincide or are disjoint.

Whether two labels overlap is the same on both sides.  If an overlap has
witnesses $\gamma,\zeta$ on the first side, extend $P'$ to match the type
of $(P,\gamma,\zeta)$.  This uses
\[
 (r+2)s\leq(m+1)s\leq M.
\]
Projection preserves the two label conditions, and
Lemma~\ref{lem:hf-atomic-transfer} preserves equality of the two values.
Hence the label sets overlap on the second side.  The converse is
symmetric.

We also need equality of the sizes of corresponding label sets.  Fix a
nonempty label.  The number $n$ of molecules of its type over $P$ equals
the number over $P'$, by Lemma~\ref{lem:hf-type-counts}.  The map
$\gamma\mapsto\phi*\gamma$ from these molecules onto the label set has
fibres of a common positive size: the stabilizer acts transitively on the
molecules, and equivariance maps one fibre bijectively to another.
Choose molecules $\gamma,\gamma'$ whose joint types with the two prefixes
match.  The fibre at $\phi*\gamma$ has size
\[
 h=\bigl|\{\zeta:
       \operatorname{tp}_M(P,\zeta)=\tau,
       \ \phi*\gamma=\phi*\zeta\}\bigr|.
\]
Partition the possible $\zeta$ by the finitely many realized types of
$(P,\gamma,\zeta)$.  On each such type the two displayed conditions are
constant, both within either structure and between structures: use
projection for the first condition and equality transfer for the second.
Each type has the same number of extensions over $(P,\gamma)$ and
$(P',\gamma')$, again using $(r+2)s\leq M$.  Summing these finite counts
shows that the fibre size $h$ agrees on the two sides.  Both value sets
therefore have size $n/h$.

Representation shows that the label sets cover the domains.  Their
nonempty labels have the same equivalence relation on both sides, where
two labels are equivalent when their value sets overlap.  For each
class, the common value set on the first side and the common value set on
the second side are finite and have equal size.  Choose a bijection
between them.  These bijections combine to a bijection $f:D_0\to D_1$,
because different classes give disjoint value sets.

There are only countably many labels: the base type sets are finite and
the forms are finite expressions over them.  No effective choice is
being assumed here.  For example, fix external enumerations of the two
countable HF universes and of the labels; take the first label of each
class and match its finite value sets in their external enumeration
orders.  These orders belong only to the proof.

If $f$ pairs $x$ with $x'$, a representative label $(\phi,\tau)$ of their
class gives molecules $\gamma,\gamma'$ with
$x=\phi*\gamma$, $x'=\phi*\gamma'$ and common joint type $\tau$ over
$P,P'$.  Thus the extended assignments are represented alike.  If
$x=x_j$ for a retained position, equality transfer forces $x'=x'_j$.
This proves the last assertion as well.
\end{proof}

\begin{proof}[Proof of Theorem~\ref{thm:hf-transfer}]
We prove by induction on $C^m$ formulas that represented-alike assignments
agree on their truth values.  Equality, membership, the atom predicate,
and the input relations agree by Lemma~\ref{lem:hf-atomic-transfer}; for
atomic formulas using two distinct variables, project onto their two
representing molecule blocks.  With one repeated variable, equality is
true and self-membership is false in $\mathrm{HF}$; a diagonal base
relation is false on a set and otherwise is determined by the corresponding
atom coordinate in its molecule type.  Boolean connectives preserve
agreement.

Consider a quantifier binding a variable $x$.  Retain the assigned values
of the other variable names, dropping the old value of $x$ if there is
one.  At most $m-1$ positions remain.  Lemma~\ref{lem:hf-bijection} gives
a bijection of the domains under which every extension by $x$ is again
represented alike.  The induction hypothesis therefore makes this a
bijection between the sets of values satisfying the quantified body.
Existential truth agrees, and for every natural number $h$ the statement
that this set has exactly $h$ elements agrees.  This proves the counting
step even if that set is infinite: a bijection preserves finiteness and
every finite cardinality.  Universal quantifiers follow in the same way
or by negation.  Empty assignments are represented alike, so all $C^m$
sentences have the same truth value.
\end{proof}

Theorem~\ref{thm:hf-transfer} imposes no bound on the rank of an object or
the depth of its representing form. The type-extension arguments require
$4s\leq M$ for equality transfer and $(m+1)s\leq M$ for the domain
bijections. Orbit homogeneity is used at length $2s$ for representation
and at length at most $ms$ when extending a retained assignment.
These requirements are independent of rank. In the next section, the
support theorem places every code and recursive constituent of a
polynomial computation inside $D_i$. We then construct bounded-variable
halt and output sentences evaluated on $D_i$, to which the transfer
theorem applies.
